\noindent \textbf{PI.} Andi Quinn (Assistant Professor, UC Santa Cruz), Email: aquinn1@ucsc.edu, Address: 1156 High Street, Santa Cruz CA 95060

\noindent \textbf{Grad Student Researcher.} Yusheng Zheng (Graduate Student, UC Santa Cruz), Email: yzhen165@ucsc.edu, Address: 1156 High Street, Santa Cruz CA 95060


\noindent \textbf{Introduction.}
eBPF has massive adoption on performance-critical kernel paths: it powers the Kubernetes networking datapath through Cilium~\cite{cilium}, Meta's production layer-4 load balancer Katran~\cite{katran}, fleet-wide observability and profiling through BCC and bpftrace~\cite{bcc,bpftrace}, runtime security monitoring through Tracee and KRSI~\cite{tracee,krsi_lwn}, and, most recently, CPU scheduling through sched\_ext~\cite{sched_ext_lwn}.  Because these programs execute on every packet, system call, or scheduling event, each additional nanosecond of per-event cost is multiplied billions of times across a fleet.  eBPF's safety comes from an in-kernel compilation pipeline: a userspace compiler lowers each program to bytecode, the kernel verifier proves memory safety and termination~\cite{verifier,gershuni2019simple}, and a deliberately simple just-in-time compiler (the kernel JIT) translates the verified bytecode to native code one instruction at a time~\cite{linux_bpf_jit}.

Unfortunately, eBPF users pay for this safety with performance.  Our characterization over 27 microbenchmarks extracted from production eBPF programs shows that the same C code runs up to twice as slow through the eBPF pipeline as when compiled directly to native code (1.57$\times$ geomean on x86-64).  The gap is structural.  The single-pass kernel JIT cannot recognize multi-instruction patterns: a 64-bit variable rotate---a single \texttt{ROL} instruction natively---arrives at the JIT as eight bytecode instructions and leaves as 15 machine instructions.  Fixing the JIT in place is unattractive: any change requires upstream acceptance and a long release cycle, enlarges the trusted computing base (TCB) that years of formal-methods work have hardened~\cite{nelson2020specification,vishwanathan2023verifying}, and grows per-architecture kernel code.  Consequently, existing eBPF optimizers such as K2~\cite{xu2021synthesizing}, Merlin~\cite{mao2024merlin}, and EPSO~\cite{zhu2025epso} stay in userspace and rewrite programs within the eBPF instruction set; they can neither emit native instructions that the instruction set lacks, nor exploit facts---stable map contents, biased branches---that become visible only after a program is deployed under its production workload.

We propose work that will close eBPF's native-performance gap while leaving the kernel verifier as the sole safety authority and keeping TCB growth minimal.  The proposed work attacks the gap at three points in a program's lifecycle: an open benchmark that measures where the gap comes from; verified native operations applied when a program is loaded; and workload-guided re-optimization applied after deployment.  \textbf{Every component of the proposed work already exists as a working prototype}, open sourced in our repository at \url{https://github.com/eunomia-bpf/bpf-bench}: the benchmark suite, the Kops kernel modules and compile-time recognizer, and the BpfReJIT shim and pass engine.  Kops, our load-time extension interface with Lean~4 equivalence proofs (arXiv:2606.24213~\cite{zheng2026kops}), already speeds up eBPF microbenchmarks by up to 24\% and production datapath applications by up to 12\% on stock-verifier kernels, and BpfReJIT, our deployment-time framework, transparently re-optimizes the programs of six unmodified production applications on stock kernels.  If successful, the proposed work will (1) establish techniques that make verified kernel extensions performance-competitive with native code and (2) deliver open-source infrastructure---a benchmark suite, kernel operation modules, and a userspace re-optimization framework---that the eBPF community can adopt directly.

\noindent \textbf{Related Work.}
Existing eBPF optimizers rewrite programs within the instruction set before load: K2 synthesizes shorter bytecode~\cite{xu2021synthesizing}, and Merlin~\cite{mao2024merlin} and EPSO~\cite{zhu2025epso} apply peephole and superoptimization passes.  Such tools inherit the instruction set's expressiveness ceiling and make all decisions before the deployment workload is visible.  Formal-methods work such as Jitterbug~\cite{nelson2020specification} and Agni~\cite{vishwanathan2023verifying} proves properties of the \textit{existing} JIT and verifier but does not extend the pipeline with new operations; our proposal builds on the same philosophy by proving, in Lean~4, that each new native operation is equivalent to the bytecode sequence the verifier checks.  Alternative extension mechanisms trade safety for performance: kernel modules trust the author completely, and Rust-based extension frameworks trust a language toolchain rather than an independent load-time checker~\cite{jia2025rex}.  Profile-guided optimization is standard for userspace binaries (\eg BOLT~\cite{panchenko2019bolt}) but has no counterpart for deployed eBPF programs.  The PI and her graduate student are uniquely skilled for the proposed research due to their expertise in extending eBPF to new domains~\cite{zheng2025extending,yang2025egpu}, tracing low-level software behavior~\cite{Quinn_2018_sledgehammer}, and building novel debugging tools for software~\cite{Quinn_2022_omnitable} and heterogeneous systems~\cite{Ma_2022_bnw}.

\noindent \textbf{Research Plan.}  We will construct techniques and open-source infrastructure that close the gap between verified eBPF code and native code while holding one invariant constant: every optimized program is checked by the stock kernel verifier before it runs, so the optimizer never joins the safety argument.  The work will first complete and release an optimization benchmark built from six production eBPF applications, characterizing where the pipeline loses performance.  Then, we will complete Kops~\cite{zheng2026kops}, an interface through which kernel modules add new operations to the pipeline; each operation carries a \textit{proof sequence} of vanilla eBPF instructions that the existing verifier checks on every load and a \textit{native emit} that the JIT compiles, with Lean~4 proofs that the two forms are equivalent.  Finally, we will complete BpfReJIT, a userspace framework that re-optimizes deployed programs using facts visible only at runtime, resubmitting every candidate through the stock verifier and JIT.  Our preliminary upper-bound study (2.358$\times$ on Cilium) quantifies the long-term reward of this research line.

\textit{Thrust 1: An Open Benchmark for eBPF Runtime Optimization.}
We will complete and release the first benchmark suite for eBPF runtime optimization, built from six production applications (Cilium, Katran, Tracee, Tetragon, BCC, and an eBPF continuous profiler) with 146 eBPF programs, 42 microbenchmarks, real workloads, and hidden correctness checks that detect unsound transformations.  The suite runs unmodified upstream binaries on stock kernels across x86-64 and ARM64 with a fully automated measurement methodology.  We will use the suite to deliver a characterization of where the pipeline loses performance and as the shared evaluation substrate for Thrusts 2 and 3, and we will release our LLM-agent exploration harness, which has already discovered optimizations of up to 34\% on suite programs.

\textit{Thrust 2: Verified Native Operations at Load Time.}
Next, we will complete Kops, which lets kernel modules introduce new operations without modifying the kernel core.  Our preliminary implementation provides seven hardware-idiom operations (\eg rotate, conditional select, bit extract); it speeds up microbenchmarks by up to 24\% on x86-64 and 22\% on ARM64---recovering 42\% of the gap to native code---and improves Cilium and Katran datapath throughput by up to 12\%~\cite{zheng2026kops}.  The funded work addresses the open problem our results expose: \textit{profitability policy}.  Applying every matched site can hurt (a coverage-maximizing policy regresses Katran to 0.995$\times$), so we will develop a cost-model-guided policy that decides per site whether a replacement pays off on the target workload.  We will extend the operation set toward the map- and helper-adjacent patterns that dominate real datapaths---idioms alone recover only 5.4\% of the achievable gap on Cilium---and publish the operations, Lean~4 proofs, and kernel patches through an upstream RFC.

\textit{Thrust 3: Workload-Guided Re-Optimization After Deployment.}
Finally, we will complete BpfReJIT, which re-optimizes deployed eBPF programs using deployment-time facts.  An in-process shim intercepts an unmodified application's \texttt{BPF\_PROG\_LOAD} calls and attachment operations; a pass engine rewrites raw bytecode using runtime side inputs (captured map contents, per-site PMU branch profiles); and the shim submits every candidate through the stock verifier and JIT, so a rejected candidate leaves the original program running.  Preliminary work validates the load-time path end to end on all six benchmark applications.  The funded work will answer three open questions: (1) \textit{when} runtime facts are profitable---our profile-guided branch-layout pass improved one complete six-application run but regressed a rerun, motivating stability-aware gating evaluated on held-out runs; (2) \textit{live replacement coverage} across links, perf events, tracepoints, XDP, and tail-call program arrays; and (3) \textit{phase-change recovery}---detecting when a specialization assumption expires and reverting to the general program.

\noindent \textbf{Timeline and Artifacts.} The planned research duration is 12 months.  We split the duration into four quarters and articulate the deliverables in each quarter below:

\textit{Quarter 1 (Q1).} We will complete Thrust 1 and deliver the open-source benchmark suite and gap characterization report.

\textit{Quarter 2 (Q2).} We will complete the core of Thrust 2, delivering the Kops modules, Lean~4 proofs, and profitability policy, and post an upstream RFC.

\textit{Quarter 3 (Q3).} We will start Thrust 3 and deliver the complete load-time re-optimization evaluation with stability-aware gating across all six applications.

\textit{Quarter 4 (Q4)}. We will complete Thrust 3, delivering live-replacement coverage, phase-change recovery experiments, and a final report with a feasibility study of verifier-checked native execution.

%\noindent \textbf{Funding Amount and Effort.} The estimated budget for 2021 is \$50,000 (direct cost, no overhead). This includes salary and benefits for one Ph.D. level graduate student, and one term of tuition. The graduate student will lead the research design and evaluation.
